% Раздел: Протокол оценки (глава «Метод и экспериментальная установка»)

\subsection{Протокол оценки}
\label{sec:eval_protocol}

\subsubsection*{Метрики качества}

Для количественной оценки результатов переноса стиля применяется набор из пяти
метрик, каждая из которых измеряет отдельный аспект качества генерации.

\textbf{DINO-style}~--- косинусное сходство между \texttt{[CLS]}-токенами модели
DINO ViT-B/8~\cite{caron2021dino} для сгенерированного изображения и стилевого
референса. Эта метрика является \textit{основным} критерием оценки качества переноса
стиля: авторы B-LoRA~\cite{frenkel2024blora} используют именно DINO-сходство в качестве
ключевой метрики (Table~1 оригинальной статьи), поскольку DINO-признаки чувствительны
к текстуре и художественной манере независимо от семантики сцены. Вычисляется по модели
\texttt{facebook/dino-vitb8} (Hugging Face \texttt{transformers}~\cite{wolf2020transformers})
как среднее по батчу косинусное сходство нормализованных \texttt{[CLS]}-эмбеддингов.

\textbf{CLIP-style}~--- косинусное сходство между эмбеддингами CLIP~\cite{radford2021clip}
сгенерированного изображения и стилевого референса. Метрика оценивает степень
соответствия визуального стиля эталону: чем выше значение, тем точнее воспроизведена
художественная манера. Вычисляется по модели CLIP ViT-L/14 на нормализованных векторах
с использованием L2-нормировки перед скалярным произведением; абсолютные значения CLIP-style
в задаче переноса стиля ограничены конструкцией CLIP-пространства, однако при значениях
$\gtrsim 0{,}7$ интерпретация метрики становится осторожной (см. ограничения в гл.~<<Результаты и анализ>>).

\textbf{CLIP-content}~--- косинусное сходство между эмбеддингом CLIP сгенерированного
изображения и текстовым эмбеддингом входного промпта. Метрика оценивает, насколько
хорошо сохраняется семантическое содержание при переносе стиля: высокое значение
означает, что стилизованное изображение соответствует заданному промпту.

\textbf{FID} (\textit{Fréchet Inception Distance})~\cite{heusel2017fid}~--- расстояние
Фреше между распределениями активаций сети Inception-v3 на реальных и сгенерированных
изображениях. Вычисляется с использованием библиотеки clean-fid~\cite{parmar2022cleanfid}
на выборке из датасета ArtBench-10 (500~реальных изображений соответствующего жанра против
100~сгенерированных). Низкое значение FID соответствует высокому фотореализму и близости
к целевому стилевому распределению.

\textbf{LPIPS} (\textit{Learned Perceptual Image Patch Similarity})~\cite{zhang2018lpips}~---
перцептивное сходство, вычисляемое на основе активаций сети VGG. Используется как
вспомогательная метрика для оценки текстурного соответствия сгенерированного изображения
и стилевого референса. Низкое значение LPIPS соответствует высокому перцептивному
сходству (метрика расстояния).

\begin{figure}[h]
    \centering
    \includegraphics[width=\textwidth]{eval_pipeline.png}
    \caption{Схема пайплайна оценки. Верхний ряд~--- линейная последовательность
             обучения и инференса: одно стилевое изображение DS1 используется
             для обучения LoRA-адаптера $\Theta_\text{style}$; на этапе
             инференса 100 промптов MS-COCO val2017 (DS2) пропускаются через
             FLUX.1-dev с подключённым адаптером и дают 100~стилизованных
             изображений. Нижний ряд~--- пять метрик, вычисляемых на полученной
             выборке. Стрелки сбоку показывают вспомогательные входы: стилевой
             референс участвует в~DINO-style и CLIP-style; ArtBench-10 (DS3)
             служит реальной выборкой для FID. Все генерации детерминированы
             ($\mathtt{seed}=42$); агрегация выполняется усреднением по~100~промптам и
             4~опорным изображениям внутри стиля.}
    \label{fig:eval_pipeline}
\end{figure}

\subsubsection*{Схема оценки}

Полный протокол включает на один метод
$N_\text{gen}^{(1)} = N_\text{style} \times N_\text{ref}\times N_\text{prompts}
= 2 \times 4 \times 100 = 800$~генераций (два стиля Van~Gogh и Monet, по
4~стилевых референса в каждом, 100~промптов MS-COCO val2017); для основного
сравнения трёх методов (Phase~4) суммарно $N_\text{gen}=N_\text{methods}\cdot N_\text{gen}^{(1)}=2\,400$~генераций. Дополнительно
$1\,000$~генераций для оценки FID на подмножестве ArtBench-10
($100~\text{промптов} \times 10~\text{стилей ArtBench}$, seed фиксирован = 42).

Для каждого метода:
\begin{enumerate}
    \item Обучить LoRA-адаптер $\Theta_\text{style}$ на одном стилевом изображении
          (по одному прогону на каждое из 4~изображений стиля).
    \item Сгенерировать изображения по 100~промптам из MS-COCO val2017
          с применением $\Theta_\text{style}$.
    \item Вычислить DINO-style (косинусное сходство \texttt{[CLS]}-токенов DINO ViT-B/8
          между сгенерированными изображениями и стилевым референсом) — основная метрика.
    \item Вычислить CLIP-style (сравнение с референсным изображением)
          и CLIP-content (сравнение с текстовым промптом).
    \item Вычислить FID между 1000~сгенерированными изображениями
          и 500~изображениями ArtBench-10 соответствующего жанра
          (\textit{approximate FID}, reduced protocol).
\end{enumerate}

Итоговые значения метрик усредняются по 4~стилевым изображениям внутри каждого стиля
(Van~Gogh, Monet). Все генерации производятся при фиксированном значении seed~=~42,
шагах сэмплирования~=~28 и guidance scale~=~3{,}5.

\subsubsection*{Сравниваемые методы}

В основном эксперименте (Phase~4) сравниваются четыре метода:

\begin{itemize}
    \item \textbf{B-LoRA-FLUX} (предлагаемый) — LoRA-адаптер на double-stream
          блоках FLUX.1, конфигурация которых (диапазон блоков, ранг~$r$, число шагов)
          определяется по результатам аблационного исследования (раздел~\ref{sec:ablations});
          при инференсе применяется только $\Theta_\text{style}$.
    \item \textbf{Full-LoRA-FLUX} — LoRA-адаптер на всех блоках FLUX.1 (без
          блочного разделения); служит для проверки гипотезы о необходимости
          селективного выбора блоков.
    \item \textbf{SplitFlux}~\cite{yang2025splitflux} --- метод блочного разделения
          стиля и содержания в FLUX.1, основанный на single-stream блоках; выступает
          в роли конкурирующего метода внутри архитектуры FLUX и позволяет оценить
          относительный вклад выбора типа блоков (double-stream vs.\ single-stream).
    \item \textbf{IP-Adapter-FLUX}~\cite{ye2023ipadapter}~--- метод на основе
          адаптера изображения, не требующий файнтюнинга. На момент проведения
          экспериментов в открытом доступе отсутствовала стабильная публичная
          реализация IP-Adapter, совместимая с актуальной версией FLUX.1-dev:
          в ходе работы было опробовано несколько кандидатных репозиториев,
          однако ни один из них не обеспечил воспроизводимого пайплайна с
          использованной версией модели и весами. По этой причине IP-Adapter
          не включён в итоговое сравнение, и данный пробел зафиксирован как
          ограничение работы. Повторение сравнения после появления стабильной
          реализации указано в перспективах дальнейших исследований.
\end{itemize}

Дополнительно в Phase~5 проводится кросс-архитектурное сравнение с
\textbf{B-LoRA-SDXL}~\cite{frenkel2024blora} --- оригинальным методом на базовой
модели SDXL --- для оценки прироста качества при переходе к архитектуре FLUX.1
(подробнее~--- параграф «Phase~5» ниже).

\subsubsection*{Аблационные исследования}
\label{sec:ablations}

Аблационное исследование проводится в несколько последовательных фаз на стиле Van~Gogh
(одно фиксированное изображение, seed~=~42). Фазы выполняются в строго определённом
порядке; результаты каждой фазы определяют конфигурацию следующей. Переход между
фазами регулируется контрольной точкой GATE~1, фиксирующей итоговое решение о выборе
целевых блоков и сценарии дальнейшей работы.

\paragraph{Анализ блоков FLUX.1 (Phase~0).}
Перед обучением каких-либо адаптеров проводится анализ специализации блоков методом
prompt injection (методология описана в разделе~\ref{sec:method}, уравнение~\eqref{eq:phase0_score}).
Для каждого из 19~double-stream блоков и подмножества single-stream блоков
$\{0, 10, 20, 30, 37\}$ вычисляется оценка $S_b$ в режимах style injection
и content injection на 200~парах промптов. Результаты Phase~0 используются для:
\begin{enumerate}
    \item определения приоритетных диапазонов блоков для аблации~D-A;
    \item выявления кандидатов на роль содержательных блоков для Phase~1b.
\end{enumerate}

\paragraph{Диагностическая серия D (Phase~1).}
Базовая проверка рабочей гипотезы. Включает четыре конфигурации:

\begin{itemize}
    \item \textbf{e00\_no\_lora\_baseline} — генерация без LoRA-адаптера; служит
          нижней границей по DINO-style и CLIP-style.
    \item \textbf{d01} — double-stream блоки~[0–18], $r{=}16$, 1\,000~шагов,
          $\text{lr}{=}5{\times}10^{-5}$, промпт \texttt{"a sks"}
          (минималистичный аналог \texttt{"A [v]"} из оригинальной статьи~\cite{frenkel2024blora}).
    \item \textbf{d02} — те же блоки и ранг, 2\,000~шагов (проверка влияния
          числа итераций дообучения).
    \item \textbf{d03} — double-stream блоки~[0–18], $r{=}32$, 1\,000~шагов
          (проверка влияния ранга).
\end{itemize}

\paragraph{Поиск адаптера содержания (Phase~1b).}
По результатам Phase~0 на трёх перекрывающихся диапазонах double-stream блоков~---
ранних $[0\text{--}8]$ (конфигурация DC-01), поздних $[9\text{--}18]$ (DC-02) и
центральной зоне $[5\text{--}12]$, гипотетически предпочтительной по Phase~0 (DC-03),~---
обучается отдельный LoRA-адаптер ($r{=}16$, те же гиперпараметры, что в~d01).
При инференсе применяется только данный адаптер; измеряются CLIP-content и DINO-style.
Адаптер признаётся содержательным ($\Theta_\text{content}$), если CLIP-content
возрастает относительно базовой линии, а DINO-style остаётся ограниченным
(см.\ критерии в разделе~\ref{sec:method}).

\paragraph{Контрольная точка GATE~1.}
По совокупным результатам Phase~1 и Phase~1b принимается одно из трёх решений:
\begin{itemize}
    \item \textbf{Сценарий~A} --- double-stream блоки обеспечивают перенос стиля
          (DINO-style в~d01 существенно выше базовой линии) \textit{и} выделен
          $\Theta_\text{content}$: выполняются полные аблации~D-A–D-P, Phase~3,
          Phase~3.2, Phase~4, Phase~4b, Phase~5, Phase~6.
    \item \textbf{Сценарий~B} --- double-stream блоки дают перенос стиля, но
          $\Theta_\text{content}$ не выделяется: выполняются аблации~D-A–D-P,
          Phase~3, Phase~4, Phase~4b, Phase~5; Phase~3.2 и Phase~6 пропускаются.
    \item \textbf{Сценарий~C} --- double-stream блоки не дают видимого переноса
          стиля: проводится дополнительный анализ альтернативных подмножеств
          блоков; конфигурация переопределяется до начала аблаций.
\end{itemize}

\paragraph{Аблация D-A: диапазон double-stream блоков (Phase~2.1).}
При фиксированном ранге и числе шагов (из лучшей конфигурации серии~D)
проверяются четыре диапазона double-stream блоков: [0–6], [0–12], [0–18], [6–18].
Цель: установить, какой диапазон блоков обеспечивает наилучший перенос стиля
по метрике DINO-style при сохранении семантической связи с промптом (CLIP-content).

\paragraph{Аблация D-B: ранг LoRA (Phase~2.2).}
При фиксированном диапазоне блоков (лучший по D-A) проверяются значения
$r \in \{4, 16, 32, 64\}$. Значение $r{=}64$ соответствует конфигурации
оригинальной статьи B-LoRA~\cite{frenkel2024blora} и включается для прямой
сопоставимости результатов.

\paragraph{Аблация D-C: число шагов (Phase~2.3).}
При фиксированных блоках и ранге (лучшие по D-A и D-B) проверяются значения
$\text{steps} \in \{500, 1\,000, 2\,000, 4\,000\}$. Цель: определить компромисс
между качеством переноса стиля и временем дообучения.

\paragraph{Аблация D-P: текстовый промпт обучения (Phase~2.4).}
При финальной конфигурации блоков, ранга и числа шагов (лучшие по D-A–D-C)
проверяются четыре варианта обучающего промпта:
\texttt{"a sks"}, \texttt{"a sks painting"},
\texttt{"a [v]"}, \texttt{"a [v] painting in [s] style"}.
Цель: оценить влияние текстового кондиционирования при обучении на степень
разделения стиля и содержания. Метрики: DINO-style, CLIP-content, CLIP-style
на 100~генерациях (промпты из MS-COCO val2017).

\paragraph{Аблация G: коэффициент применения адаптера (Phase~3).}
При финальной конфигурации блоков, ранга, числа шагов и промпта (лучшие по D-A–D-P)
проверяются значения lora\_scale $\in \{0{,}3;\; 0{,}5;\; 0{,}7;\; 1{,}0;\; 1{,}5;\; 2{,}0\}$.
Согласно оригинальной статье~\cite{frenkel2024blora}, оптимальное значение
для переноса стиля лежит в диапазоне~$\alpha \approx 0{,}4{-}0{,}5$.

\paragraph{Раздельные коэффициенты $\alpha_\text{style}$ и $\alpha_\text{content}$ (Phase~3.2, только Сценарий~A).}
При наличии двух независимых адаптеров ($\Theta_\text{style}$ и $\Theta_\text{content}$)
исследуется влияние их совместного масштабирования. Тестируются следующие комбинации:
$\alpha_\text{style} \in \{0{,}4;\; 0{,}5;\; 0{,}7;\; 1{,}0\}$ при фиксированном
$\alpha_\text{content} = 1{,}0$ и, обратно, $\alpha_\text{content} \in \{0{,}4;\; 0{,}5;\; 0{,}7;\; 1{,}0\}$
при фиксированном $\alpha_\text{style} = 1{,}0$. Метрики: DINO-style, CLIP-content,
CLIP-style на 100~генерациях. Цель: найти баланс между силой стилизации
и сохранением семантики содержания.

По совокупным результатам аблаций D-A–D-C, D-P и~G выбирается итоговая конфигурация
с наилучшим балансом DINO-style и CLIP-content; она применяется в основном
эксперименте Phase~4 и финальной проверке на стиле Monet.

\paragraph{Анализ ограничений (Phase~4b).}
После основного сравнения методов (Phase~4) проводится систематический анализ
типичных артефактов генерации. Рассматриваются три класса ограничений:
\begin{itemize}
    \item \textbf{Color leakage} --- проникновение цветовой палитры стилевого
          референса в области изображения, где оно нежелательно. Количественная
          оценка: средняя цветовая разница $\Delta E$ (CIELAB) между генерацией
          и нейтральным baseline по 5~промптам с нейтральным фоном, 5~генераций
          на промпт.
    \item \textbf{Background leakage} --- сохранение структурных элементов фона
          стилевого референса в генерациях. Оценка: DINO-object similarity
          (косинусное сходство DINO ViT-B/8 на патчах переднего плана относительно
          фона) по 5~промптам с явно заданным объектом.
    \item \textbf{Complex scenes} --- деградация семантики при стилизации сложных
          сцен с несколькими объектами. Оценка: CLIP-content по 5~промптам
          с~$\geq 3$ объектами, 5~генераций на промпт.
\end{itemize}
Результаты Phase~4b служат основой для обсуждения границ применимости предлагаемого
метода в разделе «Заключение».

\paragraph{Кросс-архитектурное сравнение (Phase~5).}
Для оценки прироста качества, обусловленного переходом от UNet-архитектуры к MM-DiT,
проводится сравнение с оригинальным методом B-LoRA-SDXL~\cite{frenkel2024blora}.
Используются 4~изображения стиля Van~Gogh (те же, что в основном эксперименте);
метрики: DINO-style, CLIP-style, CLIP-content, FID. Поскольку SDXL оперирует
при иных разрешениях и настройках семплирования, для обеспечения сопоставимости
генерации B-LoRA-SDXL производятся с теми же текстовыми промптами и фиксированным
seed~=~42; разрешение унифицируется до $512 \times 512$ перед вычислением метрик.

\paragraph{Количественное сравнение методов (Phase~5.2).}
Итоговое сравнение методов проводится по протоколу Table~1 оригинальной
статьи~\cite{frenkel2024blora} на парном наборе DS8 (\textit{reduced protocol}):
50~пар контент--стиль из 8~объектов DS5 и 8~стилевых референсов DS1.
Для каждой пары генерируется одно изображение при фиксированном \texttt{seed=42};
основная метрика~--- среднее косинусное сходство DINO ViT-B/8 между
сгенерированным изображением и соответствующим стилевым референсом.
Поскольку пул пар отличается от оригинального~\cite{frenkel2024blora},
абсолютные значения не сопоставимы напрямую с опубликованной таблицей;
сравнение методов на идентичном наборе входных данных остаётся корректным.

\subsubsection*{Воспроизводимость}

Генерация изображений производится скриптом \texttt{scripts/eval/generate\_images.py}
с параметрами, заданными через Hydra-конфиг (\texttt{configs/eval/generate.yaml}).
Конфигурации обучения для каждого эксперимента хранятся в \texttt{configs/experiments/}.
Все метрики (DINO-style, CLIP-style, CLIP-content, FID, LPIPS) вычисляются
скриптом \texttt{scripts/eval/compute\_metrics.py} и логируются в ClearML~\cite{clearml2024}. Прогресс экспериментов отслеживается
в \texttt{experiments/plan.md}, обновляемом скриптом \texttt{scripts/update\_exp\_plan.py}.
